\section{Introduction}
% Delete the text and write your Introduction here:
%------------------------------------



Block diagrams are a key medium to display complicated relationships or processes in an easy-to-read manner. These block diagrams have become more and more prevalent, often being included in user manuals, scientific papers and business process documentation. With a wider adoption of block diagrams in the current literature to display these complex situations, further flowchart recognition development is required for document analysis and recognition \cite{Griffiths}.

Block diagrams comprise of many complex shapes, lines or arrows that add to the complexity of understanding them using Document Artificial Intelligence (Document AI). They range from simple systems with few processes to highly interconnected subsystems, joined by numerous components. This poses a challenge for Document AI, which has many applications for document classification \cite{kang_2014_convolutional}, information extraction \cite{hwang2019postocr} and visual question answering (source). As a result, those with visual impairments may struggle to understand block diagrams, due to the difficult nature of recognizing and refining the structural semantics of block diagrams \cite{mathew_2020_docvqa}. 

In this paper, we propose a Swin transformer backbone to improve performance, with the goal of improving the work of  Shreyanshu and Lee \cite{bhushan-lee-2022-block}. The core idea of the improvement is to replace the Faster RCNN backbone proposed in their architecture, more specifically within the shape prediction of block diagrams. A block diagram generator was implemented along side the possible improvements, to add further complexity and variety for the model training. 



